% GENERATED FILE. Edit canonical lessons, then run npm run generate:books.
% canonical-source-hash: fnv1a64:3fa73b87516f0b7d
% canonical-lessons: JA-C474-shiboru, JA-C474-tsubusu, JA-C474-kuwaeru, JA-C474-chigiru, JA-C474-usumeru

\chapter{Getting the Food Ready}
\label{ch:ja-getting-the-food-ready}
\hlchaptermodality{474}

\begin{chapteropening}
\textbf{By the end of this chapter:} \emph{I can say how food is prepared.}

Complete the last lesson of chapter 474: I can say how food is prepared.
\end{chapteropening}

\section[\texorpdfstring{\ja{しぼる} (shiboru)}{しぼる (shiboru)}]{\ja{しぼる} (shiboru) — to squeeze; to wring}

\label{lesson:JA-C474-shiboru}



\begin{quote}
Before the new one: say the Japanese for to swell, then the Japanese for to let cool.
\end{quote}

\subsection*{You'll want to know: \ja{しぼる}}
\textbf{\ja{しぼる}} — \emph{shiboru} — \textquotedblleft{}to squeeze; to wring\textquotedblright{}.

\begin{quote}
\textbf{\ja{みかんを} \ja{しぼって、のみました。}} — \emph{mikan o shibotte, nomimashita} — \textquotedblleft{}I squeezed a mandarin and drank the juice.\textquotedblright{}
\end{quote}

Polite \textbf{\ja{しぼります}}, \emph{shiborimasu}; a godan verb. Past (polite): \textbf{\ja{ははが} \ja{ふくを} \ja{しぼりました。}}, \emph{haha ga fuku o shiborimashita}.

\begin{cousinweb}
\textbf{\ja{しぼる}} (\emph{shiboru}) is a native Japanese word (wago), inherited rather than borrowed from Chinese; its kanji spelling was attached to it later (JMdict).
\end{cousinweb}

\begin{grammarlens}[title={What you've built}]
One more word for describing what you lived and felt.
\end{grammarlens}

\subsection*{Your turn}
\begin{itemize}
\raggedright
  \item \emph{Say it:} \emph{shiboru}
  \item \emph{Say it:} \emph{shiboru}, once more
  \item \emph{Say it:} \emph{mikan o shibotte, nomimashita} — I squeezed a mandarin and drank the juice.
  \item \emph{Recall:} say the Japanese for to swell, then the Japanese for to let cool, then say \emph{shiboru} again
\end{itemize}

\begin{tcolorbox}[breakable,colback=teal!4,colframe=teal!35!black,title={Before you move on}]
What does \ja{しぼる} mean? (\textquotedblleft{}To squeeze; to wring\textquotedblright{}.) Say it once more.
\end{tcolorbox}
\section[\texorpdfstring{\ja{つぶす} (tsubusu)}{つぶす (tsubusu)}]{\ja{つぶす} (tsubusu) — to mash; to crush}

\label{lesson:JA-C474-tsubusu}



\begin{quote}
Before the new one: say the Japanese for to let cool, then the Japanese for to squeeze.
\end{quote}

\subsection*{You'll want to know: \ja{つぶす}}
\textbf{\ja{つぶす}} — \emph{tsubusu} — \textquotedblleft{}to mash; to crush\textquotedblright{}.

\begin{quote}
\textbf{\ja{いもを} \ja{つぶして} \ja{ください。}} — \emph{imo o tsubushite kudasai} — \textquotedblleft{}Please mash the potatoes.\textquotedblright{}
\end{quote}

Polite \textbf{\ja{つぶします}}, \emph{tsubushimasu}; a godan verb. Past (polite): \textbf{\ja{こどもが} \ja{いもを} \ja{つぶしました。}}, \emph{kodomo ga imo o tsubushimashita}.

\begin{cousinweb}
\textbf{\ja{つぶす}} (\emph{tsubusu}) is a native Japanese word (wago), inherited rather than borrowed from Chinese: the transitive partner of \emph{tsubureru}, \textquotedblleft{}to be crushed\textquotedblright{}, and its kanji was attached later as a kun'yomi (JMdict, KANJIDIC2).
\end{cousinweb}

\begin{grammarlens}[title={What you've built}]
One more word for describing what you lived and felt.
\end{grammarlens}

\subsection*{Your turn}
\begin{itemize}
\raggedright
  \item \emph{Say it:} \emph{tsubusu}
  \item \emph{Say it:} \emph{tsubusu}, once more
  \item \emph{Say it:} \emph{imo o tsubushite kudasai} — Please mash the potatoes.
  \item \emph{Recall:} say the Japanese for to let cool, then the Japanese for to squeeze, then say \emph{tsubusu} again
\end{itemize}

\begin{tcolorbox}[breakable,colback=teal!4,colframe=teal!35!black,title={Before you move on}]
What does \ja{つぶす} mean? (\textquotedblleft{}To mash; to crush\textquotedblright{}.) Say it once more.
\end{tcolorbox}
\section[\texorpdfstring{\ja{くわえる} (kuwaeru)}{くわえる (kuwaeru)}]{\ja{くわえる} (kuwaeru) — to add}

\label{lesson:JA-C474-kuwaeru}



\begin{quote}
Before the new one: say the Japanese for to squeeze, then the Japanese for to mash.
\end{quote}

\subsection*{You'll want to know: \ja{くわえる}}
\textbf{\ja{くわえる}} — \emph{kuwaeru} — \textquotedblleft{}to add\textquotedblright{}.

\begin{quote}
\textbf{\ja{しおを} \ja{すこし} \ja{くわえました。}} — \emph{shio o sukoshi kuwaemashita} — \textquotedblleft{}I added a little salt.\textquotedblright{}
\end{quote}

Polite \textbf{\ja{くわえます}}, \emph{kuwaemasu}; an ichidan verb. Past (polite): \textbf{\ja{ははが} \ja{さとうを} \ja{くわえました。}}, \emph{haha ga satō o kuwaemashita}.

\begin{cousinweb}
\textbf{\ja{くわえる}} (\emph{kuwaeru}) is a native Japanese word (wago), inherited rather than borrowed from Chinese; its kanji spelling was attached to it later (JMdict).
\end{cousinweb}

\begin{grammarlens}[title={What you've built}]
One more word for describing what you lived and felt.
\end{grammarlens}

\subsection*{Your turn}
\begin{itemize}
\raggedright
  \item \emph{Say it:} \emph{kuwaeru}
  \item \emph{Say it:} \emph{kuwaeru}, once more
  \item \emph{Say it:} \emph{shio o sukoshi kuwaemashita} — I added a little salt.
  \item \emph{Recall:} say the Japanese for to squeeze, then the Japanese for to mash, then say \emph{kuwaeru} again
\end{itemize}

\begin{tcolorbox}[breakable,colback=teal!4,colframe=teal!35!black,title={Before you move on}]
What does \ja{くわえる} mean? (\textquotedblleft{}To add\textquotedblright{}.) Say it once more.
\end{tcolorbox}
\section[\texorpdfstring{\ja{ちぎる} (chigiru)}{ちぎる (chigiru)}]{\ja{ちぎる} (chigiru) — to tear into pieces}

\label{lesson:JA-C474-chigiru}



\begin{quote}
Before the new one: say the Japanese for to mash, then the Japanese for to add.
\end{quote}

\subsection*{You'll want to know: \ja{ちぎる}}
\textbf{\ja{ちぎる}} — \emph{chigiru} — \textquotedblleft{}to tear into pieces\textquotedblright{}.

\begin{quote}
\textbf{\ja{やさいを} \ja{てで} \ja{ちぎりました。}} — \emph{yasai o te de chigirimashita} — \textquotedblleft{}I tore the vegetables by hand.\textquotedblright{}
\end{quote}

Polite \textbf{\ja{ちぎります}}, \emph{chigirimasu}; a godan verb. Past (polite): \textbf{\ja{こどもが} \ja{のりを} \ja{ちぎりました。}}, \emph{kodomo ga nori o chigirimashita}.

\begin{cousinweb}
\textbf{\ja{ちぎる}} (\emph{chigiru}) is a native Japanese word (wago), inherited rather than borrowed from Chinese. JMdict marks it as usually written in kana; its kanji spelling writes \emph{chi} with the kanji for \textquotedblleft{}thousand\textquotedblright{}, whose native reading is \emph{chi}, and the rest with the kanji of \emph{kiru}, \textquotedblleft{}to cut\textquotedblright{} (JMdict, KANJIDIC2).
\end{cousinweb}

\begin{grammarlens}[title={What you've built}]
One more word for describing what you lived and felt.
\end{grammarlens}

\subsection*{Your turn}
\begin{itemize}
\raggedright
  \item \emph{Say it:} \emph{chigiru}
  \item \emph{Say it:} \emph{chigiru}, once more
  \item \emph{Say it:} \emph{yasai o te de chigirimashita} — I tore the vegetables by hand.
  \item \emph{Recall:} say the Japanese for to mash, then the Japanese for to add, then say \emph{chigiru} again
\end{itemize}

\begin{tcolorbox}[breakable,colback=teal!4,colframe=teal!35!black,title={Before you move on}]
What does \ja{ちぎる} mean? (\textquotedblleft{}To tear into pieces\textquotedblright{}.) Say it once more.
\end{tcolorbox}
\section[\texorpdfstring{\ja{うすめる} (usumeru)}{うすめる (usumeru)}]{\ja{うすめる} (usumeru) — to dilute; to water down}

\label{lesson:JA-C474-usumeru}



\begin{quote}
Before the new one: say the Japanese for to add, then the Japanese for to tear into pieces.
\end{quote}

\subsection*{You'll want to know: \ja{うすめる}}
\textbf{\ja{うすめる}} — \emph{usumeru} — \textquotedblleft{}to dilute; to water down\textquotedblright{}.

\begin{quote}
\textbf{\ja{しょうゆを} \ja{みずで} \ja{うすめて} \ja{ください。}} — \emph{shōyu o mizu de usumete kudasai} — \textquotedblleft{}Please dilute the soy sauce with water.\textquotedblright{}
\end{quote}

Polite \textbf{\ja{うすめます}}, \emph{usumemasu}; an ichidan verb. Past (polite): \textbf{\ja{ちちが} \ja{おさけを} \ja{うすめました。}}, \emph{chichi ga osake o usumemashita}.

\begin{cousinweb}
\textbf{\ja{うすめる}} (\emph{usumeru}) is a native Japanese word (wago), inherited rather than borrowed from Chinese; its kanji spelling was attached to it later (JMdict).
\end{cousinweb}

\begin{grammarlens}[title={What you've built}]
One more word for describing what you lived and felt.
\end{grammarlens}

\subsection*{Your turn}
\begin{itemize}
\raggedright
  \item \emph{Say it:} \emph{usumeru}
  \item \emph{Say it:} \emph{usumeru}, once more
  \item \emph{Say it:} \emph{shōyu o mizu de usumete kudasai} — Please dilute the soy sauce with water.
  \item \emph{Recall:} say the Japanese for to add, then the Japanese for to tear into pieces, then say \emph{usumeru} again
\end{itemize}

\begin{tcolorbox}[breakable,colback=teal!4,colframe=teal!35!black,title={Before you move on}]
What does \ja{うすめる} mean? (\textquotedblleft{}To dilute; to water down\textquotedblright{}.) Say it once more.
\end{tcolorbox}
